\chapter{File formats}\label{app:formats}

\section{\code{LatencyReport\_C<n>.csv}}

One file per core in \code{<workload>/newLogger/}: a header, one row per completed request in
completion order, then two footer rows (per-column worst case, then averages). Columns, 1-indexed
as the scripts address them:

\begin{table}[htbp]\centering\small
\begin{tabularx}{\textwidth}{@{}r>{\raggedright\arraybackslash}p{0.42\textwidth}Y@{}}
\toprule \# & column & meaning \\ \midrule
1 & \code{RequstID} & the global \code{msg\_id}; joins with Debugger traces and the raw trace \\
2 & \code{Request Address} & \\
3 & \code{Ready Cycle} & the simulation cycle at which the request could first issue; \code{Ready + CPU} is the issue cycle, so rows are time-resolvable \\
4 & \code{CPU Latency} & issue $-$ ready \\
5--12 & \code{L1 Stall}, \code{Requst Bus}, \code{L2 Stall}, \code{L2 Access}, \code{Response Bus}, \code{L2-DRAM Bus}, \code{DRAM}, \code{L1 Access Latency} & the stage columns, \cref{ch:request-end-to-end} \\
13 & \code{Total Latency} & the sum of 5--12 \\
14 & \code{Effective Latency} & the non-overlapped part of Total \\
15 & \code{Oldest Latency} & head-of-queue latency; 0 if the request was never the oldest \\
16--22 & \code{LLC Arrival State}, \code{LLC Gate}, \code{Resp Ahead}, \code{Resp Ahead Refills}, \code{Array Ahead}, \code{Array Ahead Writes}, \code{LLC Stall State} & the mechanism trackers, \cref{ch:monitoring} \\
\bottomrule
\end{tabularx}
\end{table}

\section{\code{Summary.csv}}

One row per core: \code{Worst-case L1 Stall Latency}, \code{Worst-case Requst Bus Latency},
\code{Worst-case L2 Stall Latency}, \code{Worst-case L2 Access Latency}, \code{Worst-case
Response Bus Latency}, \code{Worst-case L2-DRAM Bus Latency}, \code{Worst-case DRAM Latency},
\code{Worst-case L1 Access Latency}, \code{Worst-case Total Latency}, \code{Worst-case Effective
Latency}, \code{Average Latency} (the mean Effective), \code{Finish Cycle} (the core's last
completion, its total execution time) and \code{Worst-case Oldest Latency}, appended last so the
earlier indices stay stable. This is the machine-readable hook every sweep reads.

\section{\code{JobReport\_C<n>.csv}}

Written for a core running a periodic task (\cref{ch:tasks-and-demo}), one row per job, flushed
as each job finishes so the file can be tailed during a run:

\begin{lstlisting}[style=plain]
job,release_cycle,finish_cycle,deadline_cycle,deadline_miss,n_accesses,skipped_periods,mean_access_lat,max_access_lat
\end{lstlisting}

The last two are the job's mean and worst access latency as the core sees it, issue to response.

\section{The task program, \code{task\_C<n>.task.csv}}

Column 1 is the row kind. \code{attr} rows set \code{period}, \code{jobs} ($-1$ = free-running),
\code{ooo}, \code{deadline}, \code{start}; phase rows are the job body, replayed once per job.
Numbers accept hex or decimal; \code{\#} lines and blank lines are ignored.

\begin{lstlisting}[style=plain]
COMPUTE,cycles
READ,addr            WRITE,addr
LINEAR,base,count,stride,rd_ratio
RANDOM,base,range,count,rd_ratio
\end{lstlisting}

\section{The sweep result, \code{results/<axis>/<suite>.csv}}

\begin{lstlisting}[style=plain]
value,benchmark,status,avg,wcTotal,wcEff,wcL1stall,wcReqBus,wcL2stall,wcL2access,wcRespBus,wcDramBus,wcDRAM,wcL1access,finish,wcOldest
\end{lstlisting}

\code{status} is \code{OK}, \code{INCOMPLETE} or \code{TIMEOUT}; the metrics are the
\code{Summary.csv} worst cases over all cores.

\section{The raw event trace}

A file written under \env{OCTOPUS\_TRACE}: a magic, records in 64\,K-record chunks, a chunk
index (offset, count, cycle and id ranges) and a trailer, so a windowed read decodes only the
chunks it needs and a file without a trailer (a killed run) is still readable. Each record is
32 bytes, little-endian:

\begin{table}[htbp]\centering\small
\begin{tabularx}{\textwidth}{@{}llY@{}}
\toprule field & type & meaning \\ \midrule
\code{cycle} & u64 & global core cycle \\
\code{msg\_id} & u64 & the message id \\
\code{addr} & u64 & the message's address \\
\code{resource} & u8 & 0 CPU, 1 L1, 2 REQ\_BUS, 3 RESP\_BUS, 4 LLC, 5 MEM\_BUS, 6 DRAM, 8 SVC\_BUS, 9 ARRAY, 10 LLC\_QUEUE, 11 FSM \\
\code{phase} & u8 & 0 ENTER, 1 SERVICE, 2 EXIT; for FSM records the event id \\
\code{kind} & u8 & 1 DEMAND, 2 GETS, 3 GETM, 4 PUTM, 5 INV, 6 MEM\_READ, 7 EVICT, 8 WB\_DATA, 9 WB\_INV, 10 SUPPLY, 11 SUPPLY\_DEFERRED, 12 RESP, 13 FILL, 14 FILL\_ROLLBACK, 15 MEM\_WRITE; 0 untagged \\
\code{core} & u8 & the message's owner \\
\code{comp} & u16 & component id (L1 0--3, LLC 10, DRAM 100; bus 0 = L1--LLC, 1 = memory bus) \\
\code{flags} & u16 & ARRAY: the action that claimed the port (1 read for a response, 4 read for a write-back or supply, 6 write of a fill, write-back or supply), bit 8 cut-in, bit 9 MSHR-resident, bit 10 write-back-buffer-resident. FSM: \code{old\_state << 8 | new\_state} \\
\bottomrule
\end{tabularx}
\end{table}

At close the Logger writes \code{<file>.names} with each controller's state and event names, in
FSM CSV order, which the converter uses to resolve FSM records. \code{tools/octoviz/convert.py
<newLogger> --trace <file>} produces \code{octoviz.parquet} (the rows), \code{occupancy.parquet}
(one busy interval per resource per message) and \code{fsm.parquet} (one row per transition).
